\documentclass[aspectratio=169]{beamer}

% ---------- Codificação e idioma ----------
\usepackage[utf8]{inputenc}
\usepackage[T1]{fontenc}

% ---------- Tema ----------
\usetheme{Madrid}
\usecolortheme{default}
\setbeamertemplate{navigation symbols}{}
\setbeamertemplate{footline}[frame number]

\usepackage{booktabs}
\usepackage{array}
\usepackage{tikz}
\usepackage[table]{xcolor}

\definecolor{ufalblue}{RGB}{0,51,102}
\setbeamercolor{structure}{fg=ufalblue}
\setbeamercolor{title}{fg=white,bg=ufalblue}
\setbeamercolor{frametitle}{fg=white,bg=ufalblue}

% ---------- Dados da capa ----------
\title[Classificação de Tumores Cerebrais em RM]{Classificação de Tumores Cerebrais em Imagens de\\ Ressonância Magnética Utilizando Transferência de Aprendizado}
\subtitle{Uma análise comparativa de desempenho, eficiência computacional e explicabilidade}
\author{Pedro Brassi Luccas}
\institute[UNIFAL]{Universidade Federal de Alfenas -- UNIFAL \\ Ciência da Computação \\[4pt] \small Orientador: Professor Doutor Luiz Eduardo da Silva}
\date{2026}

\begin{document}

% ===================== CAPA =====================
\begin{frame}
  \titlepage
\end{frame}

% ===================== SUMÁRIO =====================
\begin{frame}{Sumário}
  \tableofcontents
\end{frame}

% ===================== RESUMO =====================
\section{Resumo}
\begin{frame}{Resumo da proposta}
  \begin{itemize}
    \item Classificação de tumores cerebrais em RM: \textbf{glioma, meningioma, tumor hipofisário, ausência de tumor}
    \item Transferência de aprendizado com \textbf{ResNet50, EfficientNet-B0 e MobileNetV3}
    \item \textbf{Auditoria e deduplicação} da base antes do treinamento (SHA-256 + pHash/dHash)
    \item Avaliação por acurácia, Macro F1-Score, AUC-ROC \textit{one-vs-rest} e matriz de confusão, além de parâmetros, tamanho e tempo de inferência
    \item Explicabilidade exploratória via \textbf{Grad-CAM}
    \item[\(\to\)] Objetivo: efeito da deduplicação sobre o desempenho e melhor equilíbrio entre arquiteturas
  \end{itemize}
  \vfill
  {\footnotesize Finalidade acadêmica -- não pretende validar ferramenta de diagnóstico clínico.}
\end{frame}

% ===================== PROBLEMÁTICA =====================
\section{Problemática}

\begin{frame}{Contexto e delimitação do tema}
  \begin{itemize}
    \item RM detalha tecidos cerebrais, mas a análise exige conhecimento especializado e grande volume de imagens
    \item CNNs aprendem características visuais automaticamente, porém exigem grandes bases de treino
    \item Transferência de aprendizado reduz essa demanda, adaptando modelos pré-treinados (ImageNet)
  \end{itemize}
  \vspace{4pt}
  \begin{block}{Delimitação}
    Classificação multiclasse 2D em 4 categorias, comparando ResNet50, EfficientNet-B0 e MobileNetV3, avaliando desempenho, custo computacional, efeito da deduplicação e interpretabilidade (Grad-CAM).
  \end{block}
\end{frame}

\begin{frame}{O problema: vazamento de dados}
  \begin{alertblock}{Risco metodológico}
    Duplicatas, versões modificadas da mesma imagem ou cortes do mesmo paciente distribuídos entre treino e teste podem causar \textbf{vazamento de dados} -- a acurácia isolada não descreve adequadamente um classificador médico multiclasse.
  \end{alertblock}
  \begin{block}{Problema de pesquisa}
    Como a deduplicação da base e a escolha da arquitetura de transferência de aprendizado influenciam o desempenho, a eficiência computacional e a interpretabilidade de modelos para classificação de tumores cerebrais em RM?
  \end{block}
\end{frame}

\begin{frame}{Questão principal e questões secundárias}
  \textbf{Questão principal:} qual arquitetura (ResNet50, EfficientNet-B0, MobileNetV3) apresenta o melhor equilíbrio entre desempenho, eficiência computacional e interpretabilidade?
  \vspace{6pt}
  \textbf{Questões secundárias:}
  \begin{itemize}
    \item A remoção de duplicatas/similares modifica significativamente o desempenho?
    \item Qual arquitetura apresenta o melhor Macro F1-Score?
    \item Quais classes apresentam maior frequência de confusão?
    \item Modelos compactos mantêm desempenho competitivo frente a modelos maiores?
    \item As regiões do Grad-CAM são visualmente coerentes com as áreas relevantes?
  \end{itemize}
\end{frame}

\begin{frame}{Hipóteses}
  \textbf{Hipótese principal:} a transferência de aprendizado permite desempenho satisfatório mesmo com base de tamanho limitado.
  \begin{itemize}
    \item[\textbf{H1}] Remoção de duplicatas/similares reduz métricas artificialmente elevadas
    \item[\textbf{H2}] Trade-off desempenho \(\times\) custo: EfficientNet-B0 com melhor equilíbrio geral; ResNet50 competitiva, maior custo; MobileNetV3 compacta e próxima às demais
    \item[\textbf{H3}] Aumento de dados (apenas no treino) reduz o sobreajuste
    \item[\textbf{H4}] Grad-CAM revela decisões baseadas em regiões irrelevantes ou artefatos
  \end{itemize}
\end{frame}

\begin{frame}{Lacuna de pesquisa}
  \begin{itemize}
    \item Base \textit{Brain Tumor MRI Dataset} já relatou 98,73\%--99,96\% de acurácia em estudos recentes (DISCI; GURCAN; SOYLU, 2025; RASA et al., 2024)
    \item \textbf{Yagis et al. (2021)}: divisão em nível de fatia 2D pode inflar artificialmente a acurácia entre \textbf{29 e 55 pontos percentuais}, a depender da base neurológica (OASIS, ADNI, PPMI, base local de Parkinson)
    \item \textbf{Fateh et al. (2026)}: ao reconstruir a mesma base, apontam remoção de duplicatas/quase-duplicatas como etapa necessária contra vazamento
  \end{itemize}
  \begin{alertblock}{Lacuna}
    Não há, até o momento, trabalhos que apliquem essa auditoria diretamente sobre a base original combinada às três arquiteturas propostas.
  \end{alertblock}
\end{frame}

% ===================== JUSTIFICATIVA =====================
\section{Justificativa}
\begin{frame}{Justificativa}
  \begin{columns}[T]
    \begin{column}{0.33\textwidth}
      \textbf{Teórica}
      \small Metodologia precisa considerar qualidade dos dados e risco de vazamento, não só acurácia final. Auditoria de duplicatas é o principal diferencial metodológico.
    \end{column}
    \begin{column}{0.33\textwidth}
      \textbf{Social}
      \small Diagnóstico precoce melhora sobrevida. Contribui ao apoio computacional ao diagnóstico por imagem, com potencial benefício indireto futuro.
    \end{column}
    \begin{column}{0.33\textwidth}
      \textbf{Tecnológica}
      \small Bases públicas viabilizam estudo sem coleta clínica direta; produz pipeline reprodutível (auditoria, treino, explicabilidade).
    \end{column}
  \end{columns}
\end{frame}

% ===================== OBJETIVOS =====================
\section{Objetivos}
\begin{frame}{Objetivo geral}
  \begin{block}{}
    Desenvolver e avaliar modelos de classificação de tumores cerebrais em RM utilizando transferência de aprendizado, comparando desempenho preditivo, eficiência computacional, influência da deduplicação e explicabilidade.
  \end{block}
\end{frame}

\begin{frame}{Objetivos específicos}
  \begin{itemize}
    \item Selecionar e documentar uma base pública de ressonâncias magnéticas
    \item Analisar distribuição, origem, licença, formatos e qualidade das imagens
    \item Detectar duplicatas exatas (hash criptográfico) e similares (hash perceptual)
    \item Definir divisões de treino/validação/teste sem vazamento conhecido
    \item Implementar uma CNN simples como \textit{baseline}
    \item Aplicar ResNet50, EfficientNet-B0 e MobileNetV3 com transferência de aprendizado
    \item Comparar métricas preditivas e computacionais, com e sem deduplicação
    \item Analisar erros e matrizes de confusão por classe
    \item Aplicar Grad-CAM e discutir limitações, vieses e generalização
  \end{itemize}
\end{frame}

% ===================== METODOLOGIA =====================
\section{Metodologia}
\begin{frame}{Natureza da pesquisa e base de dados}
  \textbf{Natureza:} pesquisa aplicada, quantitativa, experimental e exploratória.
  \vspace{6pt}
  \textbf{Base de dados:} \textit{Brain Tumor MRI Dataset} (Nickparvar, 2021, Kaggle) -- 4 classes (glioma, meningioma, tumor hipofisário, ausência de tumor).
  \begin{itemize}
    \item Registro de quantidade por classe, formato, resolução, licença e fontes primárias
    \item Kaggle tratado como plataforma de redistribuição (citação da origem original)
    \item Inclusão: arquivos legíveis, rótulos válidos, licença compatível
    \item Exclusão/análise à parte: corrompidos, rótulos ausentes/ambíguos, duplicatas confirmadas, origem inadequada
  \end{itemize}
\end{frame}

\begin{frame}{Auditoria, deduplicação e divisão dos dados}
  \begin{itemize}
    \item Auditoria \textbf{antes} da divisão dos dados
    \item \textbf{SHA-256}: identifica arquivos exatamente iguais
    \item \textbf{pHash/dHash} (distância de Hamming): sinaliza similaridade visual
    \item Pares registrados, agrupados e inspecionados por amostragem -- sem exclusão automática (podem ser cortes distintos de um mesmo exame)
  \end{itemize}
  \begin{alertblock}{Divisão dos dados}
    Base não disponibiliza identificadores de paciente \(\Rightarrow\) divisão \textbf{estratificada por classe}: 70\% treino / 15\% validação / 15\% teste. Duplicatas de um mesmo grupo ficam em uma única partição; teste isolado até a definição final.
  \end{alertblock}
\end{frame}

\begin{frame}{Pré-processamento, modelos e treinamento}
  \begin{itemize}
    \item Conversão de canais, redimensionamento com preservação de proporção, normalização específica do modelo
    \item Aumento de dados (treino apenas): rotações, translações, zoom, ajustes leves de brilho/contraste
    \item CNN simples como \textit{baseline}
    \item Redes pré-treinadas: extratores congelados \(\to\) nova cabeça de classificação \(\to\) \textbf{fine-tuning parcial} com taxa de aprendizado menor
    \item Condições equivalentes entre experimentos: divisão, transformações, épocas, \textit{early stopping}, seleção de checkpoint
  \end{itemize}
  \begin{block}{Desequilíbrio de classes}
    Se a auditoria revelar desequilíbrio relevante: \textit{class weighting} ou reamostragem.
  \end{block}
\end{frame}

\begin{frame}{Plano experimental}
  \begin{table}
    \centering
    \small
    \begin{tabular}{@{}cll@{}}
      \toprule
      \textbf{Exp.} & \textbf{Modelo} & \textbf{Base} \\
      \midrule
      E1 & CNN \textit{baseline} & sem deduplicação \\
      E2 & EfficientNet-B0 pré-treinada\textsuperscript{*} & sem deduplicação \\
      E3 & CNN \textit{baseline} & deduplicada \\
      E4 & ResNet50 & deduplicada \\
      E5 & EfficientNet-B0 & deduplicada \\
      E6 & MobileNetV3 & deduplicada \\
      E7 & Grad-CAM + análise de erros & melhor modelo \\
      \bottomrule
    \end{tabular}
  \end{table}
  {\footnotesize \textsuperscript{*}Representa as arquiteturas em E2 por H2 indicar seu melhor equilíbrio geral.}
  \vspace{6pt}
  \begin{itemize}
    \item \textbf{H1 (dedup.)}: medida de forma completa em E1 vs.\ E3; representativa em E2 vs.\ E5
    \item Extensão a ResNet50/MobileNetV3: conforme cronograma
  \end{itemize}
\end{frame}

\begin{frame}{Métricas e determinação do "melhor equilíbrio"}
  \textbf{Métricas preditivas:} acurácia, precisão, sensibilidade, especificidade, F1 por classe, Macro F1, Macro AUC, AUC-ROC \textit{one-vs-rest}, matriz de confusão
  \vspace{4pt}
  \textbf{Métricas computacionais:} parâmetros totais/treináveis, tamanho (MB), tempo de inferência
  \vspace{10pt}
  \begin{alertblock}{Melhor equilíbrio}
    Determinado de forma \textbf{qualitativa e discursiva}, comparando métricas e resultados do Grad-CAM -- sem fórmula de ponderação formal única.
  \end{alertblock}
\end{frame}

\begin{frame}{Explicabilidade e tecnologias}
  \textbf{Explicabilidade:} Grad-CAM (Selvaraju et al., 2017) em amostra de acertos de alta confiança e erros mais frequentes da matriz de confusão.
  \vspace{10pt}
  \begin{table}
    \centering
    \small
    \begin{tabular}{@{}ll@{}}
      \toprule
      \textbf{Categoria} & \textbf{Ferramentas possíveis} \\
      \midrule
      Linguagem e IA & Python; PyTorch ou TensorFlow/Keras; scikit-learn \\
      Imagem e dados & OpenCV; Pillow; ImageHash; NumPy; Pandas \\
      Visualização & Matplotlib; Seaborn \\
      Ambiente & Jupyter; Google Colab; Kaggle \\
      Versionamento & Git e GitHub \\
      \bottomrule
    \end{tabular}
  \end{table}
  {\footnotesize Extensões (se o cronograma permitir): testes estatísticos formais, FLOPs/MACs, validação externa.}
\end{frame}

% ===================== CRONOGRAMA =====================
\section{Cronograma}
\begin{frame}{Cronograma da pesquisa}
  {\footnotesize Indicativo para dez meses -- ajustado em conjunto com o orientador.}
  \vspace{6pt}
  \begin{table}
    \centering
    \tiny
    \begin{tabular}{@{}l*{10}{c}@{}}
      \toprule
      \textbf{Atividade} & 1&2&3&4&5&6&7&8&9&10 \\
      \midrule
      Tema e problema          & \cellcolor{green!30}C &&&&&&&&& \\
      Revisão bibliográfica    & &\cellcolor{yellow!40}P&\cellcolor{yellow!40}P&&&&&&& \\
      Base e documentação      & &\cellcolor{green!30}C&&&&&&&& \\
      Auditoria/deduplicação   & &&&&&&&&& \\
      Pré-processamento        & &&&&&&&&& \\
      Baseline                 & &&&&&&&&& \\
      Treinamento (E1--E6)     & &&&&&&&&& \\
      Ajustes                  & &&&&&&&&& \\
      Avaliação                & &&&&&&&&& \\
      Grad-CAM (E7)            & &&&&&&&&& \\
      Extensões (se houver tempo) & &&&&&&&&& \\
      Redação                  & &\cellcolor{yellow!40}P&\cellcolor{yellow!40}P&&&&&&& \\
      Revisão/apresentação     & &&&&&&&&& \\
      \bottomrule
    \end{tabular}
  \end{table}
  {\footnotesize C: concluído. P: parcialmente realizado.}
\end{frame}

% ===================== BIBLIOGRAFIA =====================
\section{Bibliografia}
\begin{frame}[allowframebreaks]{Bibliografia principal}
  \footnotesize
  \begin{itemize}
    \item DISCI, R.; GURCAN, F.; SOYLU, A. Advanced brain tumor classification in MR images using transfer learning and pre-trained deep CNN models. \textit{Cancers}, v.~17, n.~1, p.~121, 2025.
    \item FATEH, A. et al. BRISC: annotated dataset for brain tumor segmentation and classification. \textit{Scientific Data}, 2026.
    \item HE, K.; ZHANG, X.; REN, S.; SUN, J. Deep residual learning for image recognition. In: \textit{CVPR}, 2016, Las Vegas. p.~770--778.
    \item HOWARD, A. et al. Searching for MobileNetV3. In: \textit{ICCV}, 2019, Seoul. p.~1314--1324.
    \item NICKPARVAR, M. Brain Tumor MRI Dataset. \textit{Kaggle}, 2021.
    \item SELVARAJU, R. R. et al. Grad-CAM: visual explanations from deep networks via gradient-based localization. In: \textit{ICCV}, 2017, Venice. p.~618--626.
    \item TAN, M.; LE, Q. V. EfficientNet: rethinking model scaling for convolutional neural networks. In: \textit{ICML}, 36., 2019, Long Beach. p.~6105--6114.
    \item YAGIS, E. et al. Effect of data leakage in brain MRI classification using 2D convolutional neural networks. \textit{Scientific Reports}, v.~11, n.~22544, 2021.
  \end{itemize}
  \vspace{6pt}
  \textbf{Bibliografia complementar}
  \begin{itemize}
    \item RASA, S. M. et al. Brain tumor classification using fine-tuned transfer learning models on MRI images. \textit{Digital Health}, v.~10, 2024.
  \end{itemize}
\end{frame}

% ===================== ENCERRAMENTO =====================
\begin{frame}
  \centering
  \vfill
  \begin{block}{Tese}
    \centering
    \large A transferência de aprendizado permite classificar tumores cerebrais em RM com desempenho satisfatório mesmo em uma base de tamanho limitado -- mas apenas uma auditoria rigorosa de duplicatas revela se esse desempenho é real ou artificialmente inflado por vazamento de dados entre treino e teste.
  \end{block}
  \vspace{24pt}

  {\Huge Obrigado!}
  \vspace{20pt}

  {\large Perguntas?}
  \vfill
\end{frame}

\end{document}
